\chapter{Perspetivas e Papéis numa Conversa}
\label{ch:personas}

\ficcao A biblioteca pede ao chat que responda «como diretor». O rascunho tem o tom firme que lhe pediram e informa a Ana de que o livro está pronto. O cargo imaginado resolveu a voz da frase. Continua a faltar o estado que permitiria enviá-la.

\section{Pedir uma perspetiva com trabalho definido}

Uma \emph{persona} indica a perspetiva que o modelo deve adotar. Um revisor pode procurar promessas sem suporte. Um explicador pode trocar termos técnicos por palavras do quotidiano. Define a tarefa de cada um. «Revê a resposta da reserva e assinala toda a data ou promessa que não conste da ficha» dá ao revisor algo que podemos conferir.

Quando a tarefa envolve uma decisão real, confirma quem tem autoridade para a tomar. O estilo de um diretor pode ajudar a redigir a resposta. A aplicação ou a pessoa responsável confirma se o aviso pode ser enviado.

\section{Separar a regra da mensagem que chega}

Num chat, podes escrever a regra e o pedido no mesmo campo, em blocos separados. Uma aplicação pode transmitir essas partes em mensagens com papéis distintos. Nos serviços que usam nomes como \texttt{system}, \texttt{developer} ou \texttt{user}, os papéis descrevem de onde vem a mensagem e como o serviço a trata. A especificação de cada fornecedor define a sua hierarquia~\citep{modelspec2026}. Escrever «system» dentro de um campo de texto deixa a mensagem no papel que a interface já lhe atribuiu.

\begin{promptlabbox}{Separar a regra do caso}
\textbf{Objetivo e contexto:} distinguir a regra que vale para todas as reservas da mensagem de uma pessoa. Usa a ficha fictícia R-014.

\textbf{Procedimento:} em papel ou num chat, escreve a regra permanente num bloco e a mensagem com a ficha noutro. Faz uma segunda ficha com o estado «pronto para levantamento». Conserva a regra nas duas respostas. Se estiveres a desenvolver uma aplicação, usa os campos de instrução que o serviço documenta e testa o mesmo caso.

\textbf{Resultado esperado e validação:} o estado «recebido» pede espera pelo aviso. O estado «pronto» permite confirmar o levantamento. Nenhuma ficha fornece uma hora. Identifica o campo ou a frase em que assentou cada decisão.

\textbf{Falhas possíveis:} mudar a regra juntamente com a ficha, dar a uma palavra do utilizador o poder de alterar o registo ou confundir o nome de um bloco com um papel reconhecido pelo serviço.
\end{promptlabbox}

\section{Usar perspetivas para formular perguntas}

Podes pedir a um modelo que leia uma proposta com finalidades distintas. Para alargar o horário da biblioteca, uma perspetiva pode perguntar quantas pessoas a usariam. Outra pode procurar o impacto na equipa. Uma terceira pode verificar se o orçamento sustenta a alteração. As sugestões saem do mesmo sistema. Quando duas vozes concordam, verifica a informação que as levou à conclusão antes de contar a concordância como evidência.

\begin{godmodebox}{Uma ficha para cada perspetiva}
\begin{lstlisting}
Tarefa: Rever uma proposta de alargamento do horário.
Perspetiva: pessoa que chega ao bairro depois das 19h.
Procura: horários que permitem ou impedem a visita.
Evidência: indica os horários citados no documento.
Se faltar informação: escreve a pergunta por esclarecer.
\end{lstlisting}
\end{godmodebox}

\begin{promptlabbox}{Pedir duas leituras da mesma proposta}
\textbf{Objetivo e contexto:} descobrir questões que a proposta ainda precisa de responder. Escolhe uma decisão pequena e reversível.

\textbf{Procedimento:} escreve o que já sabes e pede uma leitura por quem vai usar o serviço e outra por quem o vai operar. Em cada leitura, pede duas perguntas e a frase ou o dado que as motivou. Faz também a tua lista, para poderes comparar.

\textbf{Resultado esperado e validação:} distingue perguntas novas, afirmações verificáveis e preferências. Confirma os pressupostos na fonte ou com quem conhece o caso.

\textbf{Falhas possíveis:} tratar duas personagens como dois peritos independentes ou considerar um número inventado como confirmação de uma necessidade.
\end{promptlabbox}

\section{Retomar uma conversa sem trocar o estado}

Numa conversa longa, o modelo pode receber apenas parte do histórico. Mesmo quando recebe tudo, o uso da informação pode variar com a posição, a tarefa e o modelo~\citep{liu2024}. Para retomar, regista o objetivo, a decisão tomada e o que ficou por confirmar. «Enviar um aviso quando o livro estiver pronto» descreve um passo futuro. «Aviso enviado» descreve um acontecimento que precisa de confirmação. O capítulo~\ref{ch:context-engineering} desenvolve a seleção e a atualização desse registo.

A perspetiva de revisor pode ajudar a melhorar a frase. No próximo capítulo, vamos pôr a ficha à frente do revisor e conferir o que a frase afirma.
